\documentclass[10pt]{article}
\usepackage[margin=20mm]{geometry}
\usepackage{amsmath,amssymb,mathtools,fontspec,unicode-math,xurl,hyperref}
\IfFontExistsTF{DejaVu Serif}{\setmainfont{DejaVu Serif}}{\setmainfont{Latin Modern Roman}}
\setmathfont{latinmodern-math.otf}
\setlength{\parindent}{0pt}\setlength{\parskip}{6pt}\sloppy
\title{Initial scientific review
A strict maximum principle for nonlocal minimal surfaces}
\author{Assigned initial-review executor}
\date{}
\begin{document}
\maketitle
Paper. Serena Dipierro, Ovidiu Savin and Enrico Valdinoci, arXiv:2308.01697v2, 28 November 2024, 42 pages. Frozen identifier: M25-0302. Journal mapping: DOI 10.4171/JEMS/1675. All locations below refer to the acquired preprint, not to the publisher’s version.\par
\section*{Assessment and scope}
The main result asserts that nested \(s\)-minimal sets, \(0<s<1\), which share a boundary point in a ball of minimality coincide. The argument uses a normal separation function, a geometric integral equation, compactness away from the singular set, and a Harnack inequality.\par
One Major proof defect is supported by the inspected text: the measurable-set translation estimate (3.24) is false, and it is used to obtain an error estimate in the principal geometric equation. The counterexample below refutes precisely that auxiliary assertion; it does not refute Theorem 1.1. A known perimeter estimate provides a credible, mathematically specific repair for the minimal-set application. The severity reflects the role of the invalid derivation, not a claim that the main result is false.\par
All 42 pages were read as extracted text. The normal-coordinate calculations, the cut-off and capacity argument, the barrier estimates, the blow-up selection and the final dependency chain were examined. This is not a complete independent proof of the imported regularity and Harnack theorems. The cited arXiv version of the BV estimate was also checked visually at its title, Theorem 1.1 and Definitions 1.5–1.6; its full proof is not independently certified. Exact image hashes and observations accompany the provenance. The scientific decision and adjudication remain pending.\par
\section*{P1. The translation bound for arbitrary measurable sets}
\subsection*{Location and assertion}
Equation reference: bad. On printed/PDF page 15, equation (3.24) asserts, for all \(r>0\), \(\tau\in\mathbb R^n\) and measurable \(L\), a bound of the form \[
 |B_r\cap((L+\tau)\setminus L)|\leq C r^{n-1}|\tau|.\] It is substituted into (3.23), together with (3.17), in the last paragraph of the proof of Lemma 3.2. The resulting \(O(w(x_0)\eta)\) term enters Theorem 3.4, equation (3.26), and subsequently equations (8.12), (8.15) and (8.17), on pages 36–38.\par
The first equality in the proposed proof identifies the one-sided gain with the absolute value of \(|B_r\cap(L+\tau)|-|B_r\cap L|\). The signed difference is the gain minus the loss; taking its absolute value does not generally recover the gain. For example, with \(L=(0,1)\), \(r=2\), and \(0<\tau<1\), the gain has length \(\tau\), whereas the difference of the two intersection lengths is zero.\par
\subsection*{A fixed-set counterexample satisfying every stated condition}
A single fixed measurable set already contradicts [bad] (\#bad), even if \(C\) were allowed to depend on that set. Work in \(\mathbb R^2\) with \(r=1\). Put \[a_j=2^{-2j},\qquad b_j=2^{-(2j-1)}=2a_j,\qquad
 A=\bigcup_{j\geq1}(a_j,b_j),\qquad L=A\times\mathbb R.\] This is an open, hence measurable, set. For any positive integer \(N\), let \(h_N=2^{-(2N+2)}\) and \(\tau_N=(h_N,0)\). For \(1\leq j\leq N\), consider \[Q_{j,N}=(b_j,b_j+h_N)\times(-1/2,1/2).\] These rectangles are pairwise disjoint and satisfy \[Q_{j,N}\subset B_1\cap((L+\tau_N)\setminus L).\] Indeed, \(a_j\geq a_N=4h_N\), so \(x_1-h_N\in(a_j,b_j)\) for \(x_1\in(b_j,b_j+h_N)\). On the other hand, the interval to the right of \(b_j\) is a gap of \(A\) (the next interval, if present, starts at \(a_{j-1}=2b_j\)), so \(x_1\notin A\). Finally, \[x_1\leq \frac12+h_N\leq\frac9{16},\qquad |x_2|<\frac12,\qquad
 \left(\frac9{16}\right)^2+\left(\frac12\right)^2
 =\frac{145}{256}<1 .\] Each rectangle has area \(h_N\). Consequently \[\frac{|B_1\cap((L+\tau_N)\setminus L)|}{|\tau_N|}
 \geq N .\] There is no finite \(C\) for all \(N\). This is an analytic, infinite-family counterexample, not an extrapolation from a numerical experiment. An independent exact-rational program checks the interval, gap, disjointness and ball-containment inequalities for representative \(N\); the argument above proves them for every \(N\).\par
\subsection*{Impact and counterevidence}
The invalid bound is used for the symmetric difference between \(E_1\) and a translate of \(E_1\) in the balls covering the bad set. The intended bound is proportional to \(w(x_0)\sum_j r_j^{n-1}=w(x_0)\eta\). After normalization of the separation function, this becomes the error that must vanish when the singular-set covering is removed. Thus it occurs in the derivation of the core limiting inequality, rather than only in incidental exposition.\par
The counterexample is not \(s\)-minimal and is not offered as a counterexample to Theorem 1.1 or to Lemma 3.2 with additional perimeter assumptions. Minimality is important counterevidence: Cinti–Serra–Valdinoci prove a universal classical-perimeter bound for stable fractional-minimal sets [CSV]. The reviewed paper itself cites their work on page 38 for a perimeter estimate. Lemma 3.3 concerns a signed weighted tail integral and uses a change of variables; the present counterexample does not invalidate that separate argument.\par
\subsection*{Repair and its limits}
For a set of locally finite classical perimeter, the BV translation inequality gives \[\int_{B_r(p)}
 |\mathbf 1_E(x-\tau)-\mathbf 1_E(x)|\,dx
 \leq |\tau|\,P(E;B_{r+|\tau|}(p)).\] If \(|\tau|\leq r\) and \(P(E;B_{2r}(p))\leq C r^{n-1}\), this yields the required bound, now under an explicit hypothesis. If \(|\tau|>r\), the volume bound alone suffices. The perimeter hypothesis is available for stable/minimal sets when the relevant enlarged balls are inside the minimality region, by the scaled version of [Retained latex: \textbackslash{}cite[Theorem 1.1]\{CSV\}]. For fixed outer radius, the blow-up minimality region eventually contains those balls.\par
A correction should therefore replace (3.24) by this qualified estimate, state the additional hypothesis wherever Lemma 3.2 and Theorem 3.4 require it, and explain the uniformity of constants over the bad-set covering and the blow-up index. Alternatively, a valid signed-integral argument with controlled covering boundaries could be supplied. The unrestricted measurable-set assertion cannot be retained. The BV route is a checked repair mechanism; this review does not claim that a fully revised proof has been adopted by the authors or that every imported result has been independently proved.\par
\section*{Other checks and limitations}
The range \(0<s<1\) makes the local quadratic remainder integrable against the hypersurface kernel of order \(1+s\). The positive-part barrier estimates in Sections 6 and 7 and the normalization in Section 8 were traced to their stated local hypotheses. The final argument requires common tangent cones along the relevant subsequences, a regular normalization point separated from the singular set, and a nonzero normalized limit. Appendix A was read with those requirements in view.\par
A potentially misleading odd-difference expression in the normal-coordinate discussion was not promoted to a separate finding: the auxiliary averaging function is even, and its evenness can supply the required cancellation. No counterexample to the main strict maximum principle was established. The weak Harnack theorem, capacity approximation, and external smooth-convergence inputs were identified but not fully reproved. No author code was run. Targeted title/erratum and arXiv searches did not locate a correction or separate supplement in the inspected results; this is a bounded search result, not proof of absence. Absence of code is not a mathematical defect.\par
\section*{Supporting checks}
The precise acquired PDF has SHA256\par
2917b4207e21d77a5e439b7f484d9775f350908a9ba9b6b1f05b504cea929c7d\par
and source-manifest SHA256\par
6d20b849787535828af44c7fb15342c84f878e70f34f293f5b6d93c9f9ceab0d.\par
The source preparation token is preserved separately in the original manifest.\par
This is an initial review, not independent acceptance. Actual compiler commands, exit codes, dependency payloads, before/after hashes and PDF-page records are retained outside this scientific narrative. Initial `EMFILE` failures and the earlier incomplete report are preserved. Rendering and viewing remain separately recorded; the prior mathematical analysis and exact-rational check were reused.\par
9 S. Dipierro, O. Savin and E. Valdinoci, A strict maximum principle for nonlocal minimal surfaces, arXiv:2308.01697v2 (2024). E. Cinti, J. Serra and E. Valdinoci, Quantitative flatness results and BV-estimates for stable nonlocal minimal surfaces, J. Differential Geom. 112 (2019), 447–504, DOI 10.4310/jdg/1563242471. Theorem 1.1 was consulted in arXiv:1602.00540v2; no version equivalence beyond that consulted statement is asserted.\par
\end{document}
